\documentclass[10pt,a4paper]{article}
\usepackage[utf8]{inputenc}
\usepackage[T1]{fontenc}
\usepackage{amsfonts}
\usepackage[french]{babel}
\frenchbsetup{StandardLists=true}
\usepackage{enumitem}
\usepackage{amssymb}
\usepackage{amsmath}
\usepackage[left=1.5cm,right=1.5cm,top=1.5cm,bottom=1.5cm]{geometry}
\usepackage[dvipsnames]{xcolor}

\usepackage[T1]{fontenc}
\usepackage{fouriernc}
\usepackage{graphicx}

\usepackage{setspace}
\singlespacing

\usepackage{amssymb}
\usepackage{multicol}
\usepackage{fancyhdr}
\pagestyle{fancy}
\renewcommand\headrulewidth{1pt}
\fancyhead[L]{Lycée Denis Diderot }
\fancyhead[R]{Classe de 1 Bac Pro }
\setcounter{page}{}\fancyfoot[C]{BARÈME}
\fancyfoot[R]{Année 2024-2025}
\fancyfoot[L]{Alexis RUIZ}
\usepackage{multirow,tabularx,booktabs}

\usepackage{tcolorbox}
\usepackage{pifont}

\renewcommand{\arraystretch}{1.3}
\newcolumntype{C}[1]{>{\centering\arraybackslash }b{#1}}

\begin{document}

\begin{center}
\Large{\textbf{BARÈME D'ÉVALUATION - Transporter l'énergie électrique}}\\[2mm]
\normalsize
\textbf{Total : 20 points + Bonus énigme : +1 point (note max 21/20)}
\end{center}

\hrule 
\vspace{2mm}

\section*{\ding{202} Exercice 1 : QCM \hfill \textbf{5 points}}

\textit{1 point par question - Pas de point si erreur ou absence de réponse}

\begin{tabularx}{\textwidth}{|c|X|c|}
\hline
\textbf{Question} & \textbf{Contenu} & \textbf{Points} \\
\hline
\ding{202} & Unité de la tension efficace : Volt & 1 pt \\
\hline
\ding{203} & Formule de la loi d'Ohm : $R=\dfrac{U}{I}$ & 1 pt \\
\hline
\ding{204} & Nature du transformateur si m>1 : élévateur & 1 pt \\
\hline
\ding{205} & Formule de I : $I=\sqrt{\dfrac{P_j}{R}}$ & 1 pt \\
\hline
\ding{206} & Formule de U : $U=\sqrt{P_j \times R}$ & 1 pt \\
\hline
\end{tabularx}

\vspace{3mm}

\section*{\ding{203} Exercice 2 : Transformateur de piscine \hfill \textbf{7 points}}

\begin{tabularx}{\textwidth}{|c|X|c|}
\hline
\textbf{Q} & \textbf{Éléments de réponse} & \textbf{Pts} \\
\hline
\textbf{1} & Tension alternative (0,5) + Symbole correct (0,5) & \textbf{1} \\
\hline
\textbf{2.1} & Lecture : 2,4 div (0,5) + Calcul : $2,4 \times 5$ (0,5) + $U_{max} = 12$ V (0,5) & \textbf{1,5} \\
\hline
\textbf{2.2} & Formule : $U_{eff} = \dfrac{U_{max}}{\sqrt{2}}$ (0,5) + AN (0,5) + $U_{eff} \approx 8$ V (0,5) & \textbf{1,5} \\
\hline
\textbf{2.3} & Formule : $m = \dfrac{U_{2eff}}{U_{1eff}}$ (0,5) + AN (0,5) + $m \approx 0,03$ (0,5) & \textbf{1,5} \\
\hline
\textbf{2.4} & Abaisseur (0,75) + Justification : m < 1 ou U2 < U1 (0,75) & \textbf{1,5} \\
\hline
\end{tabularx}

\vspace{3mm}

\section*{\ding{204} Exercice 3 : Transformateurs sur pylônes \hfill \textbf{8 points}}

\begin{tabularx}{\textwidth}{|c|X|c|}
\hline
\textbf{Partie} & \textbf{Éléments de réponse} & \textbf{Pts} \\
\hline
\multicolumn{3}{|c|}{\textbf{Partie \ding{202} : S'approprier}} \\
\hline
& Tension dans les câbles : 20 000 V (ou 63 000 V) (0,5) + Tension sortie : 230 V (0,5) & \textbf{1} \\
\hline
\multicolumn{3}{|c|}{\textbf{Partie \ding{203} : Analyser - Raisonner}} \\
\hline
& Montage 2 (0,5) + Justification : générateur en tension alternative (0,5) & \textbf{1} \\
\hline
\multicolumn{3}{|c|}{\textbf{Partie \ding{204} : Réaliser}} \\
\hline
\textbf{3.1} & Lecture voie 1 : $U_{1max} = 582$ V (0,75) + voie 2 : $U_{2max} = 41$ V (0,75) & \textbf{1,5} \\
\hline
\textbf{3.2} & Formule $U_{eff} = \dfrac{U_{max}}{\sqrt{2}}$ (0,5) + $U_{1eff} = 411,6$ V (0,75) + $U_{2eff} = 29,0$ V (0,75) & \textbf{2} \\
\hline
\textbf{3.3} & Formule : $m = \dfrac{U_{2eff}}{U_{1eff}}$ (0,5) + AN (0,5) + $m \approx 0,070$ (0,5) & \textbf{1,5} \\
\hline
\ding{205} & Abaisseur (0,5) + Justification : m < 1 ou U2 < U1 (0,5) & \textbf{1} \\
\hline
\end{tabularx}

\vspace{3mm}

\section*{\ding{205} BONUS : Énigme \hfill \textbf{+1 point}}

\begin{tcolorbox}[colback=green!5!white,colframe=green!75!black]
\textbf{Réponse : 7 personnes (4 garçons et 3 filles)} ~~~ \textbf{Barème :} Réponse juste avec démonstration : +1 pt • Réponse juste seule : +0,5 pt
\end{tcolorbox}

\vspace{2mm}

\begin{tcolorbox}[colback=orange!5!white,colframe=orange!75!black,title=\textbf{Conseils de correction}]
\textbf{Unités :} -0,25 pt par unité manquante/incorrecte • \textbf{Calculs :} Valoriser la démarche même si résultat faux • \textbf{Justifications :} Accepter si idée principale présente • \textbf{Arrondis :} Accepter résultats proches si méthode correcte
\end{tcolorbox}

\end{document}